\documentclass[11pt]{article}
\usepackage[a4paper, total={5.5in, 10in}]{geometry}
\usepackage{hyperref} % for url link
\setlength{\parskip}{0.25cm}
\usepackage{graphicx}
\usepackage{wrapfig}
\usepackage{natbib}
\setlength{\parindent}{0in}
\setlength{\bibsep}{0.0pt}
\usepackage{pdfpages}
\usepackage{amsmath}
\usepackage{tikz}
%\renewcommand{\arraystretch}{1.7}

\begin{document}

\section{Deriving a Holling Type-II functional response}

We will begin with the simplest model of prey capture as a function of prey density. If we define the total number of prey captured, $y$ (units [D], for density or mass per unit area), as a function of the amount of time searched, $T_s$ (units [T] for time), the prey density, $x$ (units [D]), and the rate of searching, $a$ (units [T$^{-1}$]), then we can write:
    \begin{equation}
    y = axT_s. \label{eq:type1}
    \end{equation}

This equation is a linear functional response; in other words, a Type-I functional response. For many predator-prey systems, this sufficiently describes the rates at which predators remove prey from populations. One of the assumptions this makes is that time is  spent searching and not handing prey, whether that is preparing (e.g., killing, breaking into bite-sized pieces, removing husks) or digesting prey. For some predator-prey systems handling prey takes a sufficient amount of time that we want to include it in our model. We do this by writing that by writing that the total time, T$_t$ (units [T]) is made up of two parts: the time spent searching defined above and the time spent handling prey times, $b$ (units [T][D$^{-1}$]), the number of prey defined above:
    \begin{equation}
    T_t = T_s + by.
    \end{equation}
We can rearrange the equation to isolate $T_s$ as $ T_s =  T_t - by$, and substitute it into eqn.~\ref{eq:type1}:
    \begin{equation}
    y = ax\left(T_t - by\right).
    \end{equation}
If we now solve for $y$, which is, again, what we want to model, then we can do the following algebraic rearrangement:
\begin{equation}
\begin{subequations}
\begin{aligned}
    y &= aT_tx - abyx \\
    y + abyx &= aT_tx \\
    y\left(1 - abx\right) &= aT_tx \\
     y &= \frac{aT_tx}{1 + abx}. \label{eq:type2}
\end{aligned}
\end{subequations}
\end{equation}

We now want to turn this into a \emph{rate} to go into our system of differential equations. The units of our differential equations, d$N/$d$t$,  are [D][T$^{-1}$], so the right-hand side needs to be the same to balance them. This means we want the right-hand terms to be [D][T$^{-1}$] as well. To turn eqn.~\ref{eq:type2} into units of [D][T$^{-1}$], we need to divide both sides by the total time of the experiment, $T_t$, and change the symbol for prey density, $x$, to match our text as $V$. Now, we can plug it in as a term in our differential equation:

\begin{equation}
\begin{aligned}
    \frac{\mathrm{d}V}{\mathrm{d}t} &= \text{prey growth} - \left(\frac{aV}{1 + abV}\right)P \\
    \frac{\mathrm{d}P}{\mathrm{d}t} &=  h\left(\left(\frac{aV}{1 + abV}\right)P\right) - \text{predator death}.
\end{aligned}
\end{equation}
Here, $h$ is a function that represents the conversion of prey biomass to predator biomass.

\section{Table with parameters, units, and dimensions}
    \begin{tabular}{l l l}
    \hline
    Parameter/variable & Units & Dimension \\
    \hline
    $y$ & [D] & prey density removed \\
    $x$ & [D] & prey density present \\
    $T_t$ & [T] & time \\
    $T_s$ & [T] & time \\
    $N$ & [D] & density \\
    $V$ & [D] & density \\
    $P$ & [D] & density \\
    $a$ & [T$^{-1}$] & predator search rate \\
    $b$ & [T][D$^{-1}$] & handling time per unit density of prey\\
    \hline
    \end{tabular}

\section{Functional response graphs}
We can plot the two functional responses, linear and saturating, as prey removal by predators as a function of prey density.

    \begin{tikzpicture}
      \draw[->] (-0.5, 0) -- (5, 0) node[right] {Prey density~$(V)$};
      \draw[->] (0, -0.5) -- (0, 5) node[above] {Prey removal (rate)};
      \draw[scale=1, domain=-0:5, smooth, variable=\x, black] plot ({\x}, {4*\x/(0.5 + \x)});
      \draw [black] (0,0) -- (5,5);
      \draw (0, 2) node [left] {$\frac{a}{2}$};
      \draw (0.5, 0) node [below] {$h^{-1}$};
      \draw [black, dashed] (0,2) -- (0.5,2);
      \draw [black, dashed] (0.5,0) -- (0.5,2);
      \draw (0, 4) node [left] {$a$};
      \draw [black, dashed] (0,4) -- (5,4);
      \draw (5, 5) node [right] {Type I};
      \draw (5, 3.33) node [right] {Type II};
    \end{tikzpicture}

    \end{document}
